% Generated draft: requires mathematical and source review.
\documentclass[11pt]{article}
\usepackage[margin=25mm]{geometry}
\usepackage{fontspec}
\setmainfont{Latin Modern Roman}
\usepackage{amsmath,amssymb,mathtools,unicode-math}
\setmathfont{Latin Modern Math}
\usepackage{enumitem,needspace,xurl,hyperref}
\setlength{\parindent}{0pt}
\setlength{\parskip}{5pt}
\setlength{\emergencystretch}{3em}
\newcommand{\N}{\mathbb{N}}
\newcommand{\Z}{\mathbb{Z}}
\newcommand{\Q}{\mathbb{Q}}
\newcommand{\R}{\mathbb{R}}
\newcommand{\C}{\mathbb{C}}
\newcommand{\F}{\mathbb{F}}
\newcommand{\A}{\mathbb{A}}
\newcommand{\PP}{\mathbb{P}}
\newcommand{\E}{\mathbb{E}}
\newcommand{\eps}{\varepsilon}
\newcommand{\dd}{\,\mathrm{d}}
\begin{document}
\section*{UnsolvedMath ID: 20002520}
\textbf{Source problem code:} catalogue entry
\subsection{Definitions and mathematical statement}
The supplied source-verification account identifies this extraction as Answers 1 and 2 to a workshop question about generalizing compact Sobolev embeddings to complete open manifolds, rather than as a self-contained open problem. The original antecedent is not fully reproduced. A precise scalar framework for the discussed compactness is a smooth complete noncompact Riemannian manifold \((M^n,g)\) without boundary, with the nonnegative self-adjoint scalar Laplacian \(\Delta=-\operatorname{div}\nabla\) on \(L^2(M)\), and
\[
 H^1(M)=\{u\in L^2(M):\nabla u\in L^2(M)\},\qquad
 \|u\|_{H^1}^2=\int_M(|u|^2+|\nabla u|^2)\,\dd\mathrm{vol}_g.
\]
Here the essential spectrum \(\sigma_{\mathrm{ess}}(\Delta)\) excludes isolated eigenvalues of finite multiplicity. In this form-domain setting, compactness of \(H^1(M)\hookrightarrow L^2(M)\) is equivalent to compact resolvent, equivalently \(\sigma_{\mathrm{ess}}(\Delta)=\varnothing\). The extracted answers do not specify the full family of Sobolev exponents or weighted spaces intended in the preceding question.

Carron's answer says that geometric hypotheses are necessary and bounded Ricci curvature is insufficient. The source attributes to Donnelly [3] the implication
\[
 \operatorname{Ric}_g\ge-(n-1)g
 \quad\Longrightarrow\quad
 \sigma_{\mathrm{ess}}(\Delta)\cap[0,(n-1)^2/4]\ne\varnothing.
\]
The squared threshold is explicitly restored in the supplied verification account. It also attributes to Donnelly and Li [4] that a Cartan--Hadamard manifold, meaning complete, simply connected, and of nonpositive sectional curvature, has empty scalar essential spectrum if its sectional curvatures tend uniformly to \(-\infty\) at infinity. A precise reading of the last hypothesis is that for every \(A>0\) there is compact \(K_A\subset M\) such that every tangent 2-plane at every \(x\notin K_A\) has curvature below \(-A\). These are reported source results, not newly established assertions.

Answer 2 says the discussion can be generalized to weighted Sobolev spaces with weights constrained by geometry at infinity, and points to Joyce's book [7] for ALE manifolds. Neither weights, norms, asymptotic orders, nor an embedding theorem are given; no arbitrary-weight compactness claim follows.

\textbf{Illustrative family in the supplied aid.} For \(n\ge2\) and \(\alpha>0\), take a smooth complete manifold which outside a compact set has one end
\[
 (R,\infty)\times S^{n-1},\qquad
 g=\dd r^2+f_\alpha(r)^2g_{S^{n-1}},\qquad f_\alpha(r)=e^{r^\alpha},
\]
with the round unit sphere metric and a smooth compact interior. The aid reports
\[
 \sigma_{\mathrm{ess}}(\Delta)=
 \begin{cases}
 [0,\infty),&0<\alpha<1,\\
 [(n-1)^2/4,\infty),&\alpha=1,\\
 \varnothing,&\alpha>1,
 \end{cases}
\]
and hence compactness of the inhomogeneous embedding \(H^1\hookrightarrow L^2\) exactly when \(\alpha>1\). This concerns essential spectrum; additional isolated eigenvalues created by the compact interior are not excluded.

For precision, putting \(q=n-1\), radial conjugation on a spherical harmonic of eigenvalue \(\lambda_\ell\ge0\) gives the potential
\[
 V_\ell(r)=\frac{q^2\alpha^2}{4}r^{2\alpha-2}
 +\frac{q\alpha(\alpha-1)}2r^{\alpha-2}
 +\lambda_\ell e^{-2r^\alpha}.
\]
The radial sectional curvature is
\[
 K_{\mathrm{rad}}=-\alpha^2r^{2\alpha-2}
                  -\alpha(\alpha-1)r^{\alpha-2},
\]
and, when \(n\ge3\), tangent sectional curvatures are
\[
 K_{\mathrm{tan}}=e^{-2r^\alpha}-\alpha^2r^{2\alpha-2}.
\]
Thus the aid's model has curvature tending to zero for \(0<\alpha<1\), to \(-1\) for \(\alpha=1\), and to \(-\infty\) for \(\alpha>1\). The dimension restriction is essential: in dimension one the displayed warping does not produce the stated transition. This example illustrates the answers, without furnishing the missing general weighted-Sobolev problem.
\subsection{Short English statement}
This is contextual workshop discussion of Rellich compactness on complete open manifolds. It records geometric obstructions, a curvature-at-infinity criterion, and a reference to weighted ALE spaces. The supplied aid adds an explicit warped-end family whose essential spectrum and compactness change at \(\alpha=1\). The missing preceding question prevents interpreting the extraction as a new open conjecture.
\subsection{Sources}
\begin{itemize}
\item Local UnsolvedMath registry, permanent ID 20002520, catalogue entry; canonical \texttt{aim-physics-notes.json}, notes[27].
\item \emph{Questions arising in open problem sessions in AIM workshop on \(L^2\)-harmonic forms in geometry and string theory}, notes by Anda Degeratu and Mark Haskins, 17 March 2004, page 1, Answers 1 and 2; first answer attributed to Carron. Supplied URL: \url{https://aimath.org/WWN/l2harmonic/l2harmonic.pdf}. Identification and threshold correction are reported by the supplied verification, not independently checked here.
\item Original reference labels: Donnelly [3], Donnelly and Li [4], Joyce's book [7] for ALE manifolds. Full bibliographic entries and the exact theorem hypotheses are not supplied.
\item Machine-generated research attempt 1: context-only/not-a-problem classification, reported rigor proved and low novelty confidence. Its warped-end contribution is not peer reviewed and requires independent verification before citation.
\end{itemize}
\end{document}
